\documentclass[a4paper, 12pt]{article}

% Essential Packages
\usepackage{amsmath}
\usepackage{amssymb}
\usepackage[top=25mm, bottom=25mm, left=25mm, right=25mm]{geometry}
\usepackage{pgfplots}
\usepackage[inline]{enumitem}
\usepackage{mathtools}
\usepackage{tikz}
\usepackage{fancyhdr}
\usepackage{setspace}
\usepackage{hyperref}

% Misc
\pgfplotsset{compat=1.18}
\usepgfplotslibrary{polar}
\usepgfplotslibrary{fillbetween}
\setlength{\parskip}{3pt}
\setstretch{1.15}
\renewcommand{\headrulewidth}{0pt}
\renewcommand{\footrulewidth}{0pt}
\fancyhf{}

\begin{document}
\pagestyle{empty}

\begin{center}
\textbf{2025-2026 Fall \\MAT123 Final\\09/01/2026}
\end{center}

\begin{enumerate}[leftmargin=*,label=\textbf{\arabic*.}]

\item \begin{enumerate}[leftmargin=*,label=\textbf{\alph*.}]
\item Two vertices of a triangle lie on the curve $y=\sqrt{4-x^2}$ and one vertex is at the origin. One edge is parallel to the $x$-axis. What is the maximum area of the triangle?

\item \begin{enumerate}[leftmargin=*,label=\textbf{\roman*.}]
\item Evaluate $\displaystyle\lim_{x\to\frac{\pi}2}\frac{\displaystyle\int_x^{\pi/2}\ln(\sin t)\,dt}{\sin x-1}$.
\item Evaluate the sum $\displaystyle\lim_{n\to\infty}\sum_{i=1}^n\frac3n\sqrt{2+\frac{3i}{n}}$. (\textit{Hint: Riemann sum})
\end{enumerate}
\end{enumerate}

\item \begin{enumerate}[leftmargin=*,label=\textbf{\alph*.}]

\item Set up the integral (but do not evaluate) that gives the area under the line $\sqrt7\cdot y=x+1$, left to the line $x=\frac34$, and inside the circle $x^2+y^2=1$.

\item \begin{enumerate}[leftmargin=*,label=\textbf{\roman*.}]
\item Evaluate $\displaystyle\int\frac{\cos^3x+\cos x}{\sin^3x-\sin^2x}\,dx$.
\item Evaluate $\displaystyle\int_{-\pi}^{\pi}\left(1+\sin^3x\cdot\sqrt{9x^2+4\cos^2x}\right)\,dx$.
\end{enumerate}
\end{enumerate}

\item \begin{enumerate}[leftmargin=*,label=\textbf{\alph*.}]
\item Determine if the improper integral $\displaystyle\int_0^3\frac{(1+e^{-x})(x^3+x+1)}{x^5}\,dx$ converges or diverges with the Comparison Test. You cannot use the Limit Comparison Test.

\item \begin{enumerate}[leftmargin=*,label=\textbf{\roman*.}]
\item The region $R$ is bounded by $y=x^3+x,\: y=2,$ and $x=0$. Find the volume of the object if $R$ is rotated around the $x$-axis.
\item The region $R$ is bounded by $y=x^3+x,\: y=2,$ and $x=0$. Find the volume of the object if $R$ is rotated around the $y$-axis.
\end{enumerate}
\end{enumerate}

\item \begin{enumerate}[leftmargin=*,label=\textbf{\alph*.}]
\item Determine the radius and interval of convergence of $\displaystyle\sum_{n=0}^\infty\frac{(-1)^n(x-3)^n}{2^n\cdot n}$.
\item \begin{enumerate}[leftmargin=*,label=\textbf{\roman*.}]
\item Is $\displaystyle\sum_{n=2}^\infty\frac1{n(\ln n)^2}$ convergent or divergent?
\item Is $\displaystyle\sum_{n=1}^\infty\frac1{3^{n\sin\left(\frac1n\right)}}$ convergent or divergent?

\end{enumerate}
\end{enumerate}
\end{enumerate}

\newpage
\pagestyle{fancy}
\renewcommand{\headrulewidth}{1pt}
\fancyhead[R]{\textit{2025-2026 Fall Final}}
\fancyhead[L]{\textit{Hacettepe MAT123 Exam Solutions Version 2}}

%Last update: 03/09/2025 ??:??

\begin{enumerate}[leftmargin=*,label=\textbf{\arabic*.}]

\item \begin{enumerate}[leftmargin=*,label=\textbf{\alph*.}]
\item Let $0 \le x \le 2$. The base of the triangle parallel to the $x$-axis has length $2x$, and its height is $y = \sqrt{4-x^2}$. The area function is
\[A(x)=\frac{2x\cdot y}{2}=x\sqrt{4-x^2}.\]

Differentiating with respect to $x$,
\[\frac{dA}{dx}=\sqrt{4-x^2}+x\cdot\frac{-2x}{2\sqrt{4-x^2}}=\frac{4-2x^2}{\sqrt{4-x^2}}.\]

Setting $\frac{dA}{dx}=0$ gives $4-2x^2=0\implies x=\sqrt{2}$.

Evaluating the area at $x = \sqrt{2}$,
\[A\left(\sqrt{2}\right)=\sqrt{2}\sqrt{4-\left(\sqrt2\right)^2}=\sqrt2\cdot\sqrt{2}=\boxed2.\]

\item \begin{enumerate}[leftmargin=*,label=\textbf{\roman*.}]
\item Evaluating the limit using L'Hôpital's Rule (form $\left[\frac{0}{0}\right]$):
\[\lim_{x \to \frac{\pi}{2}} \frac{\int_{x}^{\pi/2} \ln(\sin t) \, dt}{\sin x - 1} = \lim_{x \to \frac{\pi}{2}} \frac{-\ln(\sin x)}{\cos x}\]

Applying L'Hôpital's Rule again (form $\left[\frac{0}{0}\right]$):
\[=\lim_{x\to\frac{\pi}{2}}\frac{-\frac{\cos x}{\sin x}}{-\sin x}=\lim_{x\to\frac{\pi}{2}} \frac{\cos x}{\sin^2 x}=\frac{0}{1}=\boxed0.\]

\item Using the definite integral definition of a Riemann sum:
\[\lim_{n \to \infty} \sum_{i=1}^{n} \frac{b-a}{n} f\left(a + i \frac{b-a}{n}\right) = \int_{a}^{b} f(x) \, dx\]

Here, $a=2$, $b=5$, $\frac{b-a}{n}=\frac{3}{n}$, and $f(x)=\sqrt{x}$. Therefore:
\[\lim_{n \to \infty}\sum_{i=1}^{n}\frac{3}{n}\sqrt{2+\frac{3i}{n}}=\int_{2}^{5}\sqrt{x}\,dx=\left[\frac{2}{3}x^{3/2}\right]_{2}^{5}=\boxed{\frac{2}{3}\left(5\sqrt5-2\sqrt2\right)}\]
\end{enumerate}
\end{enumerate}

\item \begin{enumerate}[leftmargin=*,label=\textbf{\alph*.}]

\item To find the intersection points of the line $\sqrt{7}y = x+1$ and the circle $x^2 + y^2 = 1$:
\[y^2 = \frac{(x+1)^2}{7} \implies x^2 + \frac{x^2 + 2x + 1}{7} = 1 \implies 8x^2 + 2x - 6 = 0 \implies 4x^2 + x - 3 = 0\]
Solving $(4x-3)(x+1) = 0$ gives $x = -1$ and $x = \frac{3}{4}$.

The area under the line and inside the circle is given by:
\[A = \int_{-1}^{3/4} \left( \frac{x+1}{\sqrt{7}} - \left(-\sqrt{1-x^2}\right) \right) dx =\boxed{ \int_{-1}^{3/4} \left( \frac{x+1}{\sqrt{7}} + \sqrt{1-x^2} \right) dx}\]

\item \begin{enumerate}[leftmargin=*,label=\textbf{\roman*.}]
\item Let $I = \int \frac{\cos^3 x + \cos x}{\sin^3 x - \sin^2 x} \, dx = \int \frac{(\cos^2 x + 1)\cos x}{\sin^2 x (\sin x - 1)} \, dx$. Substitute $u = \sin x$, so \linebreak$du = \cos x \, dx$ and $\cos^2 x = 1 - u^2$.
\[I = \int \frac{2 - u^2}{u^2(u-1)} \, du.\]

Using partial fraction decomposition:
\[\frac{2 - u^2}{u^2(u-1)} = \frac{A}{u} + \frac{B}{u^2} + \frac{C}{u-1} \implies A u (u-1) + B(u-1) + C u^2 = 2 - u^2\]

Setting $u = 0 \implies -B = 2 \implies B = -2$.

Setting $u = 1 \implies C = 1$.

Comparing $u^2$ coefficients: $A + C = -1 \implies A + 1 = -1 \implies A = -2$.

Integrating each term:
\[I = \int \left( \frac{-2}{u} - \frac{2}{u^2} + \frac{1}{u-1} \right) du = -2\ln|u| + \frac{2}{u} + \ln|u-1| + C\]

Substituting back $u = \sin x$:
\[\boxed{I = -2\ln\left|\sin x\right| + \frac{2}{\sin x} + \ln\left|\sin x - 1\right| + C}\]

\item Split the integral into two parts:
\[\int_{-\pi}^{\pi} \left(1 + \sin^3 x \sqrt{9x^2 + 4\cos^2 x}\right) dx = \int_{-\pi}^{\pi} dx + \int_{-\pi}^{\pi} \sin^3 x \sqrt{9x^2 + 4\cos^2 x} \, dx\]

The integrand $f(x) = \sin^3 x \sqrt{9x^2 + 4\cos^2 x}$ is an odd function because \linebreak$f(-x) = -\sin^3 x \sqrt{9(-x)^2 + 4\cos^2(-x)} = -f(x)$. Since the interval $[-\pi, \pi]$ is symmetric about $0$, the integral of the odd function vanishes.
\[= [x]_{-\pi}^{\pi} + 0 = \pi - (-\pi) = \boxed{2\pi}\]

\end{enumerate}
\end{enumerate}

\item \begin{enumerate}[leftmargin=*,label=\textbf{\alph*.}]
\item For $x \in (0, 3]$, we have $e^{-x} > 0 \implies 1 + e^{-x} > 1$, and $x^3 + x + 1 > 1$. Therefore,
\[\frac{(1+e^{-x})(x^3+x+1)}{x^5} \ge \frac{1}{x^5}.\]

Consider the integral $I = \int_{0}^{3} \frac{1}{x^5} \, dx$.
\[I=\int_{0}^{3} x^{-5} \, dx = \lim_{R \to 0^+} \left[ \frac{x^{-4}}{-4} \right]_{R}^{3} = \lim_{R \to 0^+} \left( -\frac{1}{4(3^4)} + \frac{1}{4R^4} \right) = \infty\]

Since $I$ diverges and $\frac{(1+e^{-x})(x^3+x+1)}{x^5} \ge \frac{1}{x^5} > 0$ on $(0, 3]$, by the Direct Comparison Test, the integral $\int_{0}^{3} \frac{(1+e^{-x})(x^3+x+1)}{x^5} \, dx$ also $\boxed{\text{diverges}}$.

\item \begin{enumerate}[leftmargin=*,label=\textbf{\roman*.}]
\item Rotating the region $R$ bounded by $y = x^3+x$, $y = 2$, and $x = 0$ around the $x$-axis using the Washer Method ($x \in [0, 1]$):
\[V_1 = \pi \int_{0}^{1} \left( (2)^2 - (x^3+x)^2 \right) dx = \pi \int_{0}^{1} \left( 4 - (x^6 + 2x^4 + x^2) \right) dx\]
\[V_1=\pi\left[4x-\frac{x^7}7-\frac{2x^5}5-\frac{x^3}3\right]_0^1=\pi\left(4-\frac17-\frac25-\frac13\right)=\boxed{\frac{328\pi}{105}}\]

\item Rotating the same region around the $y$-axis using the Shell Method.
\[V_2=2\pi\int_0^1x\left(2-(x^3+x) \right) dx = 2\pi \int_{0}^{1} (2x - x^4 - x^2) \, dx\]
\[V_2=2\pi\left[x^2-\frac{x^5}5-\frac{x^3}3\right]_0^1=2\pi\left(1-\frac15-\frac13\right)=2\pi \left(\frac7{15}\right)=\boxed{\frac{14\pi}{15}}\]
\end{enumerate}
\end{enumerate}

\item \begin{enumerate}[leftmargin=*,label=\textbf{\alph*.}]
\item Using the Ratio Test for $\sum_{n=1}^{\infty} \frac{(-1)^n (x-3)^n}{2^n \cdot n}$:
\[L = \lim_{n \to \infty} \left| \frac{a_{n+1}}{a_n} \right| = \lim_{n \to \infty} \left| \frac{(x-3)^{n+1}}{2^{n+1}(n+1)} \cdot \frac{2^n \cdot n}{(x-3)^n} \right| = \frac{|x-3|}{2} \lim_{n \to \infty} \frac{n}{n+1} = \frac{|x-3|}{2}\]
For convergence, $L < 1 \implies \frac{|x-3|}{2} < 1 \implies |x-3| < 2$.
Thus, the radius of convergence is $\boxed{R = 2}$.

Solving $|x-3| < 2$ yields $1 < x < 5$.

At $x = 1$: $\sum_{n=1}^{\infty} \frac{(-1)^n (-2)^n}{2^n \cdot n} = \sum_{n=1}^{\infty} \frac{1}{n}$, which is the harmonic series and diverges.

At $x = 5$: $\sum_{n=1}^{\infty} \frac{(-1)^n 2^n}{2^n \cdot n} = \sum_{n=1}^{\infty} \frac{(-1)^n}{n}$, which converges by the Alternating Series Test.

Therefore, the interval of convergence is $\boxed{(1, 5]}$.

\item \begin{enumerate}[leftmargin=*,label=\textbf{\roman*.}]
\item Applying the Integral Test for $f(x) = \frac{1}{x(\ln x)^2}$ on $[2, \infty)$:
\[\int_{2}^{\infty} \frac{1}{x(\ln x)^2} \, dx = \lim_{R \to \infty} \int_{\ln 2}^{\ln R} \frac{1}{u^2} \, du = \lim_{R \to \infty} \left[ -\frac{1}{u} \right]_{\ln 2}^{\ln R} = \frac{1}{\ln 2}.\]

Since the integral converges, the series $\sum_{n=2}^{\infty} \frac{1}{n(\ln n)^2}$ 1 $\boxed{\text{converges}}$.

\item Using the $n$-th Term Test for Divergence,
\[L = \lim_{n \to \infty} \frac{1}{3^{n \sin(1/n)}}\]

Let $u = 1/n$. As $n \to \infty$, $u \to 0^+$.
\[\lim_{n \to \infty} n \sin(1/n) = \lim_{u \to 0^+} \frac{\sin u}{u} = 1\implies \lim_{n \to \infty} \frac{1}{3^{n \sin(1/n)}} = \frac{1}{3^1} = \frac{1}{3} \ne 0\]

Since the limit is not 0, the series $\boxed{\text{diverges}}$ by the $n$th Term Test.
\end{enumerate}
\end{enumerate}
\end{enumerate}

\end{document}